\capitulo{v2c12}
% =====================================================================
%  Capítulo v2c12 · Volumes e áreas de superfície
% =====================================================================
\tikzset{
  v2c12 aresta/.style={draw=ebookDark,line width=.85pt,line join=round,line cap=round},
  v2c12 oculta/.style={draw=ebookDark,line width=.55pt,dash pattern=on 2.6pt off 2pt},
  v2c12 fina/.style={draw=ebookMuted,line width=.5pt,line join=round},
  v2c12 cota/.style={draw=ebookMuted,line width=.5pt,{Stealth[length=4pt]}-{Stealth[length=4pt]}},
  v2c12 rot/.style={font=\small,text=ebookText},
}

\section{Volume e capacidade: prismas e cilindros}

Quantos litros cabem em uma caixa-d'água? Quantos caminhões de concreto são
necessários para uma laje? Em quantos dias um reservatório se esvazia? Quase todo
ano o ENEM traz ao menos uma questão em que é preciso medir o \textbf{espaço
ocupado} por um objeto — e ela quase sempre vem disfarçada de situação do
cotidiano: embalagens, piscinas, silos, latas, sorvetes, velas, contêineres.

A boa notícia é que poucas ideias resolvem quase tudo: \emph{contar camadas}
(área da base $\times$ altura), \emph{o fator $\frac13$} das pirâmides e dos cones,
\emph{a fórmula da esfera} e \emph{a semelhança} entre sólidos. Neste capítulo
você vai dominar essas ideias e aprender a reconhecer quando a pergunta é sobre
volume e quando é sobre área de superfície.

\subsection{Volume é contar cubinhos em camadas}

\begin{definicao}[Volume e capacidade]
O \textbf{volume} de um sólido é a medida do espaço que ele ocupa. Medir um volume
é contar quantos \textbf{cubos unitários} cabem dentro dele — por exemplo, cubos de
1~cm de aresta, cada um com volume de $1\ \text{cm}^3$.

A \textbf{capacidade} de um recipiente é o volume do seu \emph{interior}, isto é,
quanto de líquido ou de material ele comporta. Ela costuma ser expressa em
litros (L) ou mililitros (mL).
\end{definicao}

\noindent
\begin{minipage}[c]{0.50\linewidth}
Na figura ao lado, cabem 4 cubinhos no comprimento e 3 na largura: cada
\textbf{camada} tem $4\cdot 3 = 12$ cubinhos, que é exatamente a \textbf{área da
base}. Como há 2 camadas, o bloco tem $12\cdot 2 = 24$ cubinhos.

Esse raciocínio — \emph{área da base vezes o número de camadas} — é a ideia que
sustenta quase todas as fórmulas de volume deste capítulo.
\end{minipage}\hfill
\begin{minipage}[c]{0.46\linewidth}
\centering
\begin{tikzpicture}[x={(0.72cm,0cm)},y={(0.36cm,0.26cm)},z={(0cm,0.72cm)},line join=round]
  \fill[ebookLight] (0,0,0)--(4,0,0)--(4,0,2)--(0,0,2)--cycle;
  \fill[ebookPrim!22] (0,0,2)--(4,0,2)--(4,3,2)--(0,3,2)--cycle;
  \fill[ebookPrim!42] (4,0,0)--(4,3,0)--(4,3,2)--(4,0,2)--cycle;
  \begin{scope}[draw=ebookPrim,line width=.4pt]
    \foreach \i in {1,2,3}{\draw (\i,0,0)--(\i,0,2)--(\i,3,2);}
    \draw (0,0,1)--(4,0,1)--(4,3,1);
    \foreach \j in {1,2}{\draw (0,\j,2)--(4,\j,2)--(4,\j,0);}
  \end{scope}
  \draw[v2c12 aresta] (0,0,0)--(4,0,0)--(4,3,0)--(4,3,2)--(0,3,2)--(0,0,2)--cycle
    (0,0,2)--(4,0,2)--(4,0,0) (4,0,2)--(4,3,2);
  \draw[decorate,decoration={brace,mirror,amplitude=4pt},ebookMuted]
    (0,0,-0.12)--(4,0,-0.12) node[midway,below=5pt,v2c12 rot]{$a = 4$};
  \draw[decorate,decoration={brace,mirror,amplitude=4pt},ebookMuted]
    (4.12,0,-0.05)--(4.12,3,-0.05) node[midway,right=6pt,v2c12 rot]{$b = 3$};
  \draw[decorate,decoration={brace,amplitude=4pt},ebookMuted]
    (-0.12,0,0)--(-0.12,0,2) node[midway,left=5pt,v2c12 rot]{$c = 2$};
\end{tikzpicture}
\captionof{figure}{Bloco formado por $4\cdot3\cdot2 = 24$ cubos unitários.}
\end{minipage}

\begin{formula}[Paralelepípedo reto-retângulo e cubo]
$\displaystyle V_{\text{paralelepípedo}} = a\cdot b\cdot c
\qquad\qquad V_{\text{cubo}} = a^3$
\end{formula}

\subsection{Unidades: de metro cúbico para litro}

Como $1\ \text{m} = 10\ \text{dm}$, um cubo de 1~m de aresta comporta
$10\cdot 10\cdot 10 = 1000$ cubos de 1~dm de aresta. E um cubo de 1~dm de aresta
tem capacidade de exatamente 1~litro. Daí vêm as relações mais usadas do capítulo:

\begin{formula}[Volume e capacidade]
$1\ \text{m}^3 = 1000\ \text{L}\qquad\qquad
1\ \text{dm}^3 = 1\ \text{L}\qquad\qquad
1\ \text{cm}^3 = 1\ \text{mL}$
\end{formula}

Repare: a cada “degrau” de unidade de volume (m$^3$ → dm$^3$ → cm$^3$), o número
fica \textbf{1000 vezes} maior, e não 10 vezes, porque são três dimensões.

\begin{exemplo}[Caixa-d'água de uma casa]
Uma caixa-d'água tem a forma de um paralelepípedo reto-retângulo com dimensões
internas de 1,5~m, 1,0~m e 0,8~m. Uma família de 4 pessoas consome, em média,
150~L de água por pessoa por dia. Estando a caixa cheia e sem reabastecimento,
por quantos dias ela atende a família?
\solucao
O volume interno é $V = 1{,}5\cdot 1{,}0\cdot 0{,}8 = 1{,}2\ \text{m}^3$. Como
$1\ \text{m}^3 = 1000\ \text{L}$, a capacidade é $1{,}2\cdot 1000 = 1200\ \text{L}$.
O consumo diário da família é $4\cdot 150 = 600\ \text{L}$. Logo, a caixa dura
$1200 \div 600 = 2$ dias.

\textbf{Ligação com o ENEM:} questões de reservatório quase sempre combinam
\emph{volume} com \emph{conversão para litros} e com uma \emph{taxa} (consumo por
dia, vazão por hora). Faça uma coisa de cada vez e anote as unidades.
\end{exemplo}

\subsection{O princípio de Cavalieri (sem complicação)}

Pegue uma pilha de moedas iguais e empurre-a de lado, deixando-a torta. A forma
mudou, mas o volume não mudou: continuam sendo as mesmas moedas! Essa observação
simples é o \textbf{princípio de Cavalieri}, e é ele que garante que a ideia de
“área da base vezes altura” vale para \emph{qualquer} prisma ou cilindro, reto ou
inclinado.

\begin{center}
\begin{tikzpicture}[scale=.82]
  \newcommand\vdcmoeda[2]{%
    \filldraw[fill=ebookAcc!38,draw=ebookAcc!70!black,line width=.4pt]
      (#1-1,#2) arc[start angle=180,end angle=360,x radius=1,y radius=.26]
      -- ++(0,.15) arc[start angle=0,end angle=180,x radius=1,y radius=.26] -- cycle;
    \filldraw[fill=ebookAcc!18,draw=ebookAcc!70!black,line width=.4pt]
      (#1,#2+.15) ellipse[x radius=1,y radius=.26];}
  \foreach \i in {0,...,11}{\vdcmoeda{0}{.19*\i}}
  \foreach \i/\dx in {0/0,1/.18,2/.28,3/.45,4/.62,5/.70,6/.88,7/1.05,8/1.12,9/1.30,10/1.48,11/1.58}
    {\vdcmoeda{4.2+\dx}{.19*\i}}
  \draw[ebookPrim,line width=.8pt,dash pattern=on 4pt off 2.5pt] (-1.6,1.42)--(7.6,1.42)
    node[right,font=\footnotesize,text=ebookPrim,align=left]{mesmo nível:\\seções de\\mesma área};
  \node[font=\footnotesize,text=ebookMuted] at (0,-.65) {pilha reta};
  \node[font=\footnotesize,text=ebookMuted] at (5.0,-.65) {pilha inclinada};
  \draw[v2c12 cota] (-1.35,0)--(-1.35,2.24) node[midway,left,v2c12 rot]{$h$};
\end{tikzpicture}
\captionof{figure}{Mesmas moedas, mesma altura: mesmo volume.}
\end{center}

\begin{definicao}[Princípio de Cavalieri]
Se dois sólidos têm a mesma altura e \emph{todo} plano paralelo à base os corta
em seções de \textbf{mesma área}, então eles têm o \textbf{mesmo volume}.
\end{definicao}

Assim, um prisma (ou um cilindro) cuja base tem área $A_b$ e cuja altura é $h$
tem o mesmo volume que um paralelepípedo de base com área $A_b$ e altura $h$:
cada fatia horizontal de um tem a mesma área que a fatia correspondente do
outro. Para o cilindro, a base é um círculo de área $\pi r^2$.

\begin{formula}[Prisma e cilindro (retos ou oblíquos)]
$\displaystyle V_{\text{prisma}} = A_b\cdot h
\qquad\qquad V_{\text{cilindro}} = \pi r^2 h$
\end{formula}

\begin{center}
\begin{tikzpicture}[scale=.95]
  % ---------- prisma hexagonal
  \def\R{1.15}\def\H{2.3}\def\k{.36}
  \fill[ebookLight] ({\R*cos(180)},{\k*\R*sin(180)})
    \foreach \t in {240,300,360}{ -- ({\R*cos(\t)},{\k*\R*sin(\t)}) }
    \foreach \t in {360,300,240,180}{ -- ({\R*cos(\t)},{\k*\R*sin(\t)+\H}) } -- cycle;
  \fill[ebookPrim!30] ({\R},{\H}) \foreach \t in {60,120,...,300}{ -- ({\R*cos(\t)},{\k*\R*sin(\t)+\H}) } -- cycle;
  \draw[v2c12 oculta] ({\R*cos(0)},0) \foreach \t in {60,120,180}{ -- ({\R*cos(\t)},{\k*\R*sin(\t)}) };
  \foreach \t in {60,120}{\draw[v2c12 oculta] ({\R*cos(\t)},{\k*\R*sin(\t)}) -- ({\R*cos(\t)},{\k*\R*sin(\t)+\H});}
  \draw[v2c12 aresta] ({\R*cos(180)},0) \foreach \t in {240,300,360}{ -- ({\R*cos(\t)},{\k*\R*sin(\t)}) };
  \draw[v2c12 aresta] ({\R},{\H}) \foreach \t in {60,120,...,300}{ -- ({\R*cos(\t)},{\k*\R*sin(\t)+\H}) } -- cycle;
  \foreach \t in {0,180,240,300}{\draw[v2c12 aresta] ({\R*cos(\t)},{\k*\R*sin(\t)}) -- ({\R*cos(\t)},{\k*\R*sin(\t)+\H});}
  \node[font=\small,text=ebookDark] at (0,\H) {$A_b$};
  \draw[v2c12 cota] (\R+.35,0)--(\R+.35,\H) node[midway,right,v2c12 rot]{$h$};
  \node[font=\footnotesize,text=ebookMuted] at (0,-.85) {prisma: $V = A_b\cdot h$};
  % ---------- cilindro
  \begin{scope}[xshift=4.6cm]
    \def\r{1.05}\def\ry{.36}
    \fill[ebookLight] (-\r,0) arc[start angle=180,end angle=360,x radius=\r,y radius=\ry]
      -- (\r,\H) arc[start angle=0,end angle=-180,x radius=\r,y radius=\ry] -- cycle;
    \fill[ebookPrim!30] (0,\H) ellipse[x radius=\r,y radius=\ry];
    \draw[v2c12 aresta] (0,\H) ellipse[x radius=\r,y radius=\ry];
    \draw[v2c12 aresta] (-\r,0)--(-\r,\H) (\r,0)--(\r,\H);
    \draw[v2c12 aresta] (-\r,0) arc[start angle=180,end angle=360,x radius=\r,y radius=\ry];
    \draw[v2c12 oculta] (\r,0) arc[start angle=0,end angle=180,x radius=\r,y radius=\ry];
    \fill[ebookDark] (0,\H) circle (1.1pt);
    \draw[ebookDark,line width=.6pt] (0,\H)--(\r,\H) node[midway,above,font=\small]{$r$};
    \draw[v2c12 cota] (\r+.35,0)--(\r+.35,\H) node[midway,right,v2c12 rot]{$h$};
    \node[font=\footnotesize,text=ebookMuted] at (0,-.85) {cilindro: $V = \pi r^2 h$};
  \end{scope}
  % ---------- prisma oblíquo
  \begin{scope}[xshift=9cm,x={(1cm,0cm)},y={(.42cm,.3cm)},z={(0cm,1cm)}]
    \def\s{.75}
    \fill[ebookLight] (0,0,0)--(1.5,0,0)--(1.5+\s,0,\H)--(\s,0,\H)--cycle;
    \fill[ebookLight!60!ebookPrim!40] (1.5,0,0)--(1.5,1.4,0)--(1.5+\s,1.4,\H)--(1.5+\s,0,\H)--cycle;
    \fill[ebookPrim!30] (\s,0,\H)--(1.5+\s,0,\H)--(1.5+\s,1.4,\H)--(\s,1.4,\H)--cycle;
    \draw[v2c12 oculta] (0,0,0)--(0,1.4,0)--(1.5,1.4,0) (0,1.4,0)--(\s,1.4,\H);
    \draw[v2c12 aresta] (0,0,0)--(1.5,0,0)--(1.5,1.4,0)--(1.5+\s,1.4,\H)--(\s,1.4,\H)--(\s,0,\H)--cycle
      (\s,0,\H)--(1.5+\s,0,\H)--(1.5,0,0) (1.5+\s,0,\H)--(1.5+\s,1.4,\H);
    \draw[ebookAcc,line width=.7pt,dash pattern=on 2pt off 1.5pt] (1.5+\s,0,\H)--(1.5+\s,0,0)
      node[pos=.5,right,font=\small,text=ebookAcc]{$h$};
    \draw[ebookAcc,line width=.5pt] (1.5+\s,0,0.18)--(1.5+\s-.18,0,0.18)--(1.5+\s-.18,0,0);
    \draw[ebookMuted,line width=.4pt] (1.5,0,0)--(1.5+\s,0,0);
    \node[font=\footnotesize,text=ebookMuted,align=center] at (1.1,0,-.85) {oblíquo: $V = A_b\cdot h$\\($h$ é a distância\\entre as bases)};
  \end{scope}
\end{tikzpicture}
\captionof{figure}{Prisma, cilindro e prisma oblíquo: em todos, volume = área da base $\times$ altura.}
\end{center}

\begin{exemplo}[A lata de óleo]
Uma lata de óleo tem a forma de um cilindro circular reto com 8~cm de diâmetro e
18~cm de altura (medidas internas). Usando $\pi \approx 3{,}14$, verifique se ela
comporta 900~mL de óleo.
\solucao
O raio é a \emph{metade} do diâmetro: $r = 4$~cm. Então
\[
V = \pi r^2 h \approx 3{,}14\cdot 4^2\cdot 18 = 3{,}14\cdot 288 = 904{,}32\ \text{cm}^3 .
\]
Como $1\ \text{cm}^3 = 1\ \text{mL}$, a lata comporta cerca de 904~mL, ou seja,
cabem os 900~mL, com uma pequena folga de uns 4~mL (o espaço do “ar” no topo).
Repare que o primeiro passo foi transformar diâmetro em raio — esquecer isso é
o erro mais comum em questões de cilindro.
\end{exemplo}

\begin{contexto}[Chuva medida em milímetros]
Quando a previsão do tempo informa que choveu 30~mm, isso quer dizer que a água
da chuva formaria uma lâmina de 30~mm de altura sobre uma superfície plana e
horizontal. Em $1\ \text{m}^2$ de área, isso dá
$V = 1\ \text{m}^2\cdot 0{,}03\ \text{m} = 0{,}03\ \text{m}^3 = 30\ \text{L}$.
Por isso, \textbf{1~mm de chuva corresponde a 1~litro por metro quadrado} — uma
conta de prisma que já apareceu no ENEM (índice pluviométrico).
\end{contexto}

% =====================================================================
\section{Pirâmides, cones e esferas}

Sólidos que “afinam” até um ponto, como a pirâmide e o cone, ocupam menos espaço
do que o prisma ou o cilindro de mesma base e mesma altura. Quanto menos?
Exatamente \textbf{um terço}. Se você encher de areia uma pirâmide oca e despejar o
conteúdo em um prisma oco de mesma base e mesma altura, vai precisar repetir a
operação \emph{três vezes} para encher o prisma. O mesmo acontece com o cone e o
cilindro.

\begin{center}
\begin{tikzpicture}[scale=.95]
  % ---------- pirâmide dentro do prisma
  \begin{scope}[x={(1cm,0cm)},y={(.5cm,.36cm)},z={(0cm,1cm)}]
    \def\H{2.4}
    \fill[ebookPrim!28] (1,1,\H)--(0,0,0)--(2,0,0)--cycle;
    \fill[ebookPrim!48] (1,1,\H)--(2,0,0)--(2,2,0)--cycle;
    \draw[v2c12 fina,dash pattern=on 2pt off 1.6pt] (0,0,0)--(0,2,0)--(2,2,0) (0,2,0)--(0,2,\H);
    \draw[v2c12 fina] (0,0,0)--(2,0,0)--(2,2,0)--(2,2,\H)--(0,2,\H)--(0,0,\H)--cycle
      (0,0,\H)--(2,0,\H)--(2,0,0) (2,0,\H)--(2,2,\H);
    \draw[v2c12 oculta] (1,1,\H)--(0,2,0) (0,0,0)--(0,2,0)--(2,2,0);
    \draw[v2c12 aresta] (0,0,0)--(2,0,0)--(2,2,0) (1,1,\H)--(0,0,0) (1,1,\H)--(2,0,0) (1,1,\H)--(2,2,0);
    \draw[ebookAcc,line width=.7pt,dash pattern=on 2pt off 1.5pt] (1,1,\H)--(1,1,0) node[pos=.55,right,font=\small,text=ebookAcc]{$h$};
    \fill[ebookDark] (1,1,\H) circle (1.2pt);
    \node[font=\footnotesize,text=ebookMuted,align=center] at (1.1,0,-.75) {pirâmide $=\frac13$ do prisma};
  \end{scope}
  % ---------- cone dentro do cilindro
  \begin{scope}[xshift=5.2cm]
    \def\r{1.1}\def\ry{.36}\def\H{2.6}
    \draw[v2c12 fina] (0,\H) ellipse[x radius=\r,y radius=\ry];
    \draw[v2c12 fina] (-\r,0)--(-\r,\H) (\r,0)--(\r,\H);
    \fill[ebookPrim!32] (0,\H)--(-\r,0) arc[start angle=180,end angle=360,x radius=\r,y radius=\ry] -- cycle;
    \draw[v2c12 oculta] (\r,0) arc[start angle=0,end angle=180,x radius=\r,y radius=\ry];
    \draw[v2c12 aresta] (-\r,0) arc[start angle=180,end angle=360,x radius=\r,y radius=\ry];
    \draw[v2c12 aresta] (-\r,0)--(0,\H)--(\r,0);
    \draw[ebookAcc,line width=.7pt,dash pattern=on 2pt off 1.5pt] (0,\H)--(0,0) node[pos=.5,left,font=\small,text=ebookAcc]{$h$};
    \draw[ebookAcc,line width=.7pt] (0,0)--(\r,0) node[midway,below,font=\small,text=ebookAcc]{$r$};
    \draw[ebookAcc,line width=.5pt] (0,.16)--(.16,.16)--(.16,0);
    \node[font=\small,text=ebookDark] at (.78,1.55) {$g$};
    \fill[ebookDark] (0,\H) circle (1.2pt);
    \node[font=\footnotesize,text=ebookMuted,align=center] at (0,-.8) {cone $=\frac13$ do cilindro};
  \end{scope}
\end{tikzpicture}
\captionof{figure}{Pirâmide e cone ocupam $\frac13$ do prisma e do cilindro de mesma base e mesma altura.}
\end{center}

\begin{formula}[Pirâmide e cone]
$\displaystyle V_{\text{pirâmide}} = \frac{1}{3}\,A_b\cdot h
\qquad\qquad V_{\text{cone}} = \frac{1}{3}\,\pi r^2 h$
\end{formula}

No cone circular reto, a \textbf{geratriz} $g$ (segmento do vértice até a borda
da base), a altura $h$ e o raio $r$ formam um triângulo retângulo; por isso,
$g^2 = h^2 + r^2$. Na pirâmide regular, o mesmo papel é feito pelo
\textbf{apótema da pirâmide} $m$ (altura de cada face triangular): se a base é um
quadrado de lado $\ell$, então $m^2 = h^2 + \left(\frac{\ell}{2}\right)^2$. Você
vai precisar dessas relações para calcular áreas.

\begin{exemplo}[Velas decorativas]
Uma fábrica produz dois modelos de vela maciça, ambos com 10~cm de altura: um em
forma de pirâmide quadrangular regular com base de 6~cm de lado e outro em forma
de cone circular reto com 6~cm de diâmetro da base. Usando $\pi\approx 3$, qual
modelo gasta mais parafina, e quanto a mais?
\solucao
Pirâmide: $V = \frac13\cdot 6^2\cdot 10 = \frac13\cdot 360 = 120\ \text{cm}^3$.
Cone (raio 3~cm): $V = \frac13\cdot 3\cdot 3^2\cdot 10 = 90\ \text{cm}^3$.

A pirâmide gasta $120 - 90 = 30\ \text{cm}^3$ a mais. Faz sentido: as duas têm
a mesma altura, e o círculo de diâmetro 6~cm cabe dentro do quadrado de lado
6~cm, ou seja, a base do cone é menor ($27\ \text{cm}^2$ contra $36\ \text{cm}^2$).
Com a mesma altura, os volumes ficam na mesma razão das áreas das bases.
\end{exemplo}

\subsection{A esfera}

\begin{formula}[Esfera de raio $r$]
$\displaystyle V_{\text{esfera}} = \frac{4}{3}\,\pi r^3
\qquad\qquad V_{\text{hemisfério}} = \frac{1}{2}\cdot\frac43\pi r^3 = \frac{2}{3}\,\pi r^3$
\end{formula}

Uma forma elegante de lembrar a fórmula vem de Arquimedes. Pegue um cilindro e
um cone, ambos com raio $r$ e altura $2r$, e uma esfera de raio $r$ (que cabe
exatamente dentro do cilindro). Os volumes são:

\begin{center}
\begin{tikzpicture}[scale=.95]
  \def\r{.85}\def\ry{.27}
  % cone
  \fill[ebookPrim!30] (0,2*\r)--(-\r,0) arc[start angle=180,end angle=360,x radius=\r,y radius=\ry] -- cycle;
  \draw[v2c12 oculta] (\r,0) arc[start angle=0,end angle=180,x radius=\r,y radius=\ry];
  \draw[v2c12 aresta] (-\r,0) arc[start angle=180,end angle=360,x radius=\r,y radius=\ry];
  \draw[v2c12 aresta] (-\r,0)--(0,2*\r)--(\r,0);
  \node[font=\small] at (0,-.7) {$\dfrac{2}{3}\pi r^3$};
  \node[font=\footnotesize,text=ebookMuted] at (0,-1.3) {cone};
  % esfera
  \begin{scope}[xshift=3.2cm]
    \shade[ball color=ebookPrim!35,opacity=.9] (0,\r) circle (\r);
    \draw[v2c12 aresta] (0,\r) circle (\r);
    \draw[v2c12 oculta] (\r,\r) arc[start angle=0,end angle=180,x radius=\r,y radius=\ry];
    \draw[v2c12 aresta] (-\r,\r) arc[start angle=180,end angle=360,x radius=\r,y radius=\ry];
    \fill[ebookDark] (0,\r) circle (1.1pt);
    \draw[ebookDark,line width=.6pt] (0,\r)--(\r,\r) node[midway,above,font=\small]{$r$};
    \node[font=\small] at (0,-.7) {$\dfrac{4}{3}\pi r^3$};
    \node[font=\footnotesize,text=ebookMuted] at (0,-1.3) {esfera};
  \end{scope}
  % cilindro
  \begin{scope}[xshift=6.4cm]
    \fill[ebookPrim!18] (-\r,0) arc[start angle=180,end angle=360,x radius=\r,y radius=\ry]
      -- (\r,2*\r) arc[start angle=0,end angle=-180,x radius=\r,y radius=\ry] -- cycle;
    \fill[ebookPrim!30] (0,2*\r) ellipse[x radius=\r,y radius=\ry];
    \draw[v2c12 aresta] (0,2*\r) ellipse[x radius=\r,y radius=\ry];
    \draw[v2c12 aresta] (-\r,0)--(-\r,2*\r) (\r,0)--(\r,2*\r);
    \draw[v2c12 aresta] (-\r,0) arc[start angle=180,end angle=360,x radius=\r,y radius=\ry];
    \draw[v2c12 oculta] (\r,0) arc[start angle=0,end angle=180,x radius=\r,y radius=\ry];
    \draw[v2c12 cota] (\r+.3,0)--(\r+.3,2*\r) node[midway,right,v2c12 rot]{$2r$};
    \node[font=\small] at (0,-.7) {$2\pi r^3 = \dfrac{6}{3}\pi r^3$};
    \node[font=\footnotesize,text=ebookMuted] at (0,-1.3) {cilindro};
  \end{scope}
  \node[font=\titlefont\Large,text=ebookAcc] at (10.2,.9) {1 : 2 : 3};
\end{tikzpicture}
\captionof{figure}{Cone, esfera e cilindro de mesmo raio e mesma altura $2r$: volumes na razão $1:2:3$.}
\end{center}

Ou seja: a esfera ocupa $\frac23$ do cilindro que a envolve, e o cone ocupa
$\frac13$. Esse fato é muito útil em questões de embalagens (bolas em latas,
docinhos em formas).

\begin{contexto}[O túmulo de Arquimedes]
Arquimedes (século III a.C.) considerava a relação entre a esfera e o cilindro a
sua maior descoberta e pediu que ela fosse gravada em seu túmulo. Mais de um
século depois, o político romano Cícero relatou ter encontrado, na Sicília, um
túmulo abandonado com a figura de uma esfera dentro de um cilindro — e o
identificou como o de Arquimedes.
\end{contexto}

\begin{exemplo}[Casquinha de sorvete]
Uma casquinha tem a forma de um cone circular reto com 6~cm de diâmetro na boca
e 10~cm de altura. Ela é totalmente preenchida com sorvete e, por cima, recebe uma
“bola” em forma de hemisfério, com o mesmo diâmetro da boca. Usando
$\pi\approx 3$, quantas casquinhas completas podem ser servidas com um pote de
2~L de sorvete?
\solucao
O raio é $r = 3$~cm. Parte cônica:
$V_1 = \frac13\cdot 3\cdot 3^2\cdot 10 = 90\ \text{cm}^3$.
Hemisfério: $V_2 = \frac23\cdot 3\cdot 3^3 = 54\ \text{cm}^3$.
Cada casquinha leva $90 + 54 = 144\ \text{cm}^3 = 144\ \text{mL}$.

Como $2\ \text{L} = 2000\ \text{mL}$, temos $2000\div 144 \approx 13{,}9$. Só é
possível servir \textbf{13 casquinhas completas} (a 14.ª ficaria incompleta).
Atenção ao arredondamento: aqui se arredonda para \emph{baixo}; em perguntas do
tipo “quantas viagens” ou “quantas latas comprar”, arredonda-se para \emph{cima}.
\end{exemplo}

% =====================================================================
\section{Áreas de superfície e embalagens}

Volume responde à pergunta \emph{quanto cabe}. A \textbf{área de superfície}
responde à pergunta \emph{quanto material é preciso para fabricar ou cobrir}:
papelão de uma caixa, alumínio de uma lata, tinta de um tanque, azulejos de uma
piscina, lona de uma barraca. No ENEM, as palavras \emph{pintar, revestir,
forrar, embalar, chapa, rótulo} são sinal de área — e a área é medida em
unidades quadradas (cm$^2$, m$^2$), nunca em litros.

\begin{definicao}[Área de superfície]
A \textbf{área total} $A_t$ de um sólido é a soma das áreas de todas as
superfícies que o limitam — é a área da sua \textbf{planificação}. Costumamos
separá-la em \textbf{área da base} (ou das bases) e \textbf{área lateral} $A_\ell$.
\end{definicao}

A planificação mostra de onde vêm as fórmulas dos corpos redondos. Ao
“desenrolar” a lateral de um cilindro, obtemos um \textbf{retângulo} cuja altura é
$h$ e cujo comprimento é o contorno da base, $2\pi r$. Já a lateral do cone vira
um \textbf{setor circular} de raio $g$ e arco $2\pi r$.

\begin{center}
\begin{tikzpicture}[scale=.92]
  % ---------- cilindro planificado
  \def\r{.6}\def\h{1.8}
  \pgfmathsetmacro\w{2*pi*\r}
  \fill[ebookPrim!20] (0,0) rectangle (\w,\h);
  \draw[v2c12 aresta] (0,0) rectangle (\w,\h);
  \fill[ebookPrim!35] (\w/2,\h+\r) circle (\r);
  \fill[ebookPrim!35] (\w/2,-\r) circle (\r);
  \draw[v2c12 aresta] (\w/2,\h+\r) circle (\r) (\w/2,-\r) circle (\r);
  \draw[ebookDark,line width=.6pt] (\w/2,\h+\r)--(\w/2+\r,\h+\r) node[midway,above,font=\footnotesize]{$r$};
  \draw[v2c12 cota] (0,\h+.25)--(\w/2-\r-.1,\h+.25);
  \draw[v2c12 cota] (\w/2+\r+.1,\h+.25)--(\w,\h+.25);
  \node[font=\small,fill=white,inner sep=1pt] at (\w/2,\h/2) {$A_\ell = 2\pi r\cdot h$};
  \node[font=\footnotesize,text=ebookMuted] at (\w*.82,\h+.5) {$2\pi r$};
  \draw[v2c12 cota] (\w+.3,0)--(\w+.3,\h) node[midway,right,v2c12 rot]{$h$};
  \node[font=\footnotesize,text=ebookMuted] at (\w/2,-1.65) {cilindro: retângulo $+$ 2 círculos};
  % ---------- cone planificado
  \begin{scope}[xshift=8.2cm,yshift=1.95cm]
    \def\g{2.1}\def\rb{.7}
    \pgfmathsetmacro\ang{360*\rb/\g}
    \fill[ebookPrim!20] (0,0) -- ({270-\ang/2}:\g) arc[start angle={270-\ang/2},end angle={270+\ang/2},radius=\g] -- cycle;
    \draw[v2c12 aresta] (0,0) -- ({270-\ang/2}:\g) arc[start angle={270-\ang/2},end angle={270+\ang/2},radius=\g] -- cycle;
    \fill[ebookPrim!35] (0,-\g-\rb) circle (\rb);
    \draw[v2c12 aresta] (0,-\g-\rb) circle (\rb);
    \draw[ebookDark,line width=.6pt] (0,-\g-\rb)--(\rb,-\g-\rb) node[midway,above,font=\footnotesize]{$r$};
    \node[font=\small] at ({270+\ang/2+14}:{\g*.55}) {$g$};
    \node[font=\footnotesize,text=ebookMuted] at ({270-\ang/2-12}:{\g+.15}) {arco $2\pi r$};
    \node[font=\small,fill=white,inner sep=1pt] at (0,-1.25) {$A_\ell = \pi r g$};
    \node[font=\footnotesize,text=ebookMuted] at (0,-3.6) {cone: setor $+$ círculo};
  \end{scope}
\end{tikzpicture}
\captionof{figure}{Planificações do cilindro e do cone.}
\end{center}

\begin{formula}[Áreas de superfície]
\renewcommand\arraystretch{1.45}
\begin{tabular}{lll}
\textbf{Paralelepípedo} & $A_t = 2(ab + ac + bc)$ & cubo: $A_t = 6a^2$\\
\textbf{Prisma}         & $A_t = 2A_b + A_\ell$  & $A_\ell$ = soma dos retângulos laterais\\
\textbf{Cilindro}       & $A_\ell = 2\pi r h$    & $A_t = 2\pi r^2 + 2\pi r h$\\
\textbf{Pirâmide}       & $A_t = A_b + A_\ell$   & $A_\ell$ = soma dos triângulos (altura $m$)\\
\textbf{Cone}           & $A_\ell = \pi r g$     & $A_t = \pi r^2 + \pi r g$\\
\textbf{Esfera}         & $A = 4\pi r^2$         & \\
\end{tabular}
\end{formula}

Nem sempre se usa a área total: um tanque apoiado no chão não tem o fundo
pintado; uma piscina não tem “tampa”; um telhado de pirâmide não tem base. Antes
de somar, liste \emph{quais} superfícies o problema pede.

\begin{exemplo}[Mesma capacidade, gasto diferente]
Uma empresa vai lançar um suco em embalagem de 1~L e estuda duas caixas de
papelão: um cubo de 10~cm de aresta e um paralelepípedo de
$5\ \text{cm}\times 10\ \text{cm}\times 20\ \text{cm}$. Desprezando as dobras,
qual delas gasta menos papelão?
\solucao
As duas têm a mesma capacidade:
$10^3 = 5\cdot 10\cdot 20 = 1000\ \text{cm}^3 = 1\ \text{L}$. Mas as áreas são
diferentes:
\[
A_{\text{cubo}} = 6\cdot 10^2 = 600\ \text{cm}^2,\qquad
A_{\text{caixa}} = 2(5\cdot 10 + 5\cdot 20 + 10\cdot 20) = 2\cdot 350 = 700\ \text{cm}^2 .
\]
O cubo gasta $100\ \text{cm}^2$ a menos, cerca de 14\% de economia.

Esse resultado é geral: com o volume fixo, quanto mais “compacta” a forma, menor a
área. Entre as caixas, o cubo é a mais econômica; entre as latas cilíndricas
fechadas, a mais econômica tem altura igual ao diâmetro ($h = 2r$) — uma lata
assim, de 1~L, gastaria cerca de $554\ \text{cm}^2$ de material.
\end{exemplo}

\begin{contexto}[Por que as latas não são “ótimas”?]
Uma lata de refrigerante de 350~mL tem cerca de 6,6~cm de diâmetro e 12,2~cm de
altura, bem mais alta do que larga. O motivo: a tampa e o fundo são feitos de
chapa mais grossa (e mais cara) que a lateral, a lata precisa caber na mão e ser
empilhada com eficiência. Na vida real, a “melhor” embalagem equilibra
matemática, custo e uso.
\end{contexto}

% =====================================================================
\section{Semelhança, troncos e sólidos compostos}

\subsection{Sólidos semelhantes: k, k² e k³}

Dois sólidos são \textbf{semelhantes} quando um é uma ampliação ou redução do
outro: \emph{todas} as medidas lineares (arestas, alturas, raios, diâmetros)
ficam multiplicadas pelo mesmo número $k$, chamado \textbf{razão de semelhança}.
Esferas são sempre semelhantes entre si; cubos também.

\noindent
\begin{minipage}[c]{0.50\linewidth}
Veja o que acontece ao dobrar a aresta de um cubo ($k = 2$): cada face passa a
ter $2\cdot 2 = 4$ quadradinhos (área $\times 4$) e o cubo passa a ter
$2\cdot 2\cdot 2 = 8$ cubinhos (volume $\times 8$). Isso acontece porque a área
envolve \emph{duas} medidas lineares e o volume envolve \emph{três}.
\end{minipage}\hfill
\begin{minipage}[c]{0.46\linewidth}
\centering
\begin{tikzpicture}[x={(.62cm,0cm)},y={(.3cm,.22cm)},z={(0cm,.62cm)},line join=round]
  \def\cubo#1#2{%
    \fill[ebookLight] (#1,0,0)--(#1+#2,0,0)--(#1+#2,0,#2)--(#1,0,#2)--cycle;
    \fill[ebookPrim!22] (#1,0,#2)--(#1+#2,0,#2)--(#1+#2,#2,#2)--(#1,#2,#2)--cycle;
    \fill[ebookPrim!42] (#1+#2,0,0)--(#1+#2,#2,0)--(#1+#2,#2,#2)--(#1+#2,0,#2)--cycle;}
  \cubo{0}{1}
  \draw[v2c12 aresta] (0,0,0)--(1,0,0)--(1,1,0)--(1,1,1)--(0,1,1)--(0,0,1)--cycle (0,0,1)--(1,0,1)--(1,0,0) (1,0,1)--(1,1,1);
  \node[font=\footnotesize] at (.5,0,-.55) {aresta $a$};
  \cubo{2.2}{2}
  \begin{scope}[draw=ebookPrim,line width=.4pt]
    \draw (3.2,0,0)--(3.2,0,2)--(3.2,2,2) (2.2,0,1)--(4.2,0,1)--(4.2,2,1) (2.2,1,2)--(4.2,1,2)--(4.2,1,0);
  \end{scope}
  \draw[v2c12 aresta] (2.2,0,0)--(4.2,0,0)--(4.2,2,0)--(4.2,2,2)--(2.2,2,2)--(2.2,0,2)--cycle (2.2,0,2)--(4.2,0,2)--(4.2,0,0) (4.2,0,2)--(4.2,2,2);
  \node[font=\footnotesize] at (3.2,0,-.55) {aresta $2a$};
\end{tikzpicture}
\captionof{figure}{$k=2$: área $\times 4$, volume $\times 8$.}
\end{minipage}

\begin{formula}[Razão de semelhança $k$]
$\displaystyle \frac{\text{comprimentos}}{\text{comprimentos}} = k
\qquad \frac{\text{áreas}}{\text{áreas}} = k^2
\qquad \frac{\text{volumes}}{\text{volumes}} = k^3$
\end{formula}

Como massa, capacidade e quantidade de material maciço são proporcionais ao
volume, todas elas também variam com $k^3$. Já tinta, papel de embrulho,
cobertura fina e rótulo variam com $k^2$.

\begin{exemplo}[Miniatura de perfume]
Um perfume é vendido em frasco de 100~mL. A fábrica vai produzir miniaturas
\emph{semelhantes} ao frasco original, com todas as medidas reduzidas a 40\% das
originais. Qual é a capacidade da miniatura? E, se o rótulo original tem área de
$30\ \text{cm}^2$, qual será a área do rótulo da miniatura?
\solucao
A razão de semelhança é $k = 0{,}4$. Então:
\begin{itemize}[nosep,leftmargin=1.2em]
  \item capacidade: $100\cdot k^3 = 100\cdot 0{,}064 = 6{,}4\ \text{mL}$;
  \item rótulo: $30\cdot k^2 = 30\cdot 0{,}16 = 4{,}8\ \text{cm}^2$.
\end{itemize}
Note como a capacidade “despenca”: reduzir as medidas a 40\% deixa só 6,4\% do
volume. Quem responde 40~mL (usando $k$) ou 16~mL (usando $k^2$) cai nas
armadilhas típicas desse tipo de questão.
\end{exemplo}

\subsection{Troncos de pirâmide e de cone}

Ao cortar uma pirâmide ou um cone por um plano \emph{paralelo à base}, obtemos, em
cima, uma pirâmide (ou um cone) menor, \textbf{semelhante} à original, e, embaixo,
um \textbf{tronco}. Baldes, vasos, copos, canecas e abajures costumam ter a forma
de troncos de cone. O caminho mais seguro para o volume é:
\[
V_{\text{tronco}} = V_{\text{sólido grande}} - V_{\text{sólido pequeno}},
\]
usando a semelhança para descobrir a altura que falta.

\noindent
\begin{minipage}[c]{0.40\linewidth}
\centering
\begin{tikzpicture}[scale=.5]
  \def\ry{.3}
  \fill[ebookPrim!25] (-1.5*2,4.8) arc[start angle=180,end angle=360,x radius=3,y radius=.45]
    -- (2,0) arc[start angle=0,end angle=-180,x radius=2,y radius=\ry] -- cycle;
  \fill[ebookPrim!25] (-1.5*2,4.8) -- (-2,0) arc[start angle=180,end angle=360,x radius=2,y radius=\ry] -- (3,4.8)
    arc[start angle=360,end angle=180,x radius=3,y radius=.45] -- cycle;
  \fill[ebookPrim!40] (0,4.8) ellipse[x radius=3,y radius=.45];
  \draw[v2c12 aresta] (0,4.8) ellipse[x radius=3,y radius=.45];
  \draw[v2c12 aresta] (-3,4.8)--(-2,0) (3,4.8)--(2,0);
  \draw[v2c12 aresta] (-2,0) arc[start angle=180,end angle=360,x radius=2,y radius=\ry];
  \draw[v2c12 oculta] (2,0) arc[start angle=0,end angle=180,x radius=2,y radius=\ry];
  \draw[ebookMuted,line width=.6pt,dash pattern=on 3pt off 2pt] (-2,0)--(0,-9.6)--(2,0);
  \draw[ebookMuted,line width=.5pt,dash pattern=on 1.5pt off 1.5pt] (0,4.8)--(0,-9.6);
  \fill[ebookDark] (0,-9.6) circle (2.4pt);
  \draw[ebookDark,line width=.6pt] (0,4.8)--(3,4.8) node[midway,above,font=\footnotesize]{$R = 15$};
  \draw[ebookDark,line width=.6pt] (0,0)--(2,0) node[midway,below,font=\footnotesize]{$r = 10$};
  \draw[v2c12 cota] (3.7,0)--(3.7,4.8) node[midway,right,v2c12 rot]{$24$};
  \draw[v2c12 cota] (3.7,-9.6)--(3.7,0) node[midway,right,v2c12 rot]{$x$};
\end{tikzpicture}
\captionof{figure}{Balde (tronco de cone) completado até o vértice. Medidas em cm.}
\end{minipage}\hfill
\begin{minipage}[c]{0.56\linewidth}
\begin{exemplo}[O balde]
Um balde tem a forma de um tronco de cone com 20~cm de diâmetro no fundo,
30~cm de diâmetro na boca e 24~cm de altura. Usando $\pi\approx 3$, qual é a sua
capacidade?
\solucao
Complete o tronco até formar um cone (figura). O cone pequeno, de raio 10, é
semelhante ao cone grande, de raio 15, com $k = \frac{10}{15} = \frac23$. Se $x$ é
a altura do cone pequeno:
\[
\frac{x}{x + 24} = \frac23 \;\Rightarrow\; 3x = 2x + 48 \;\Rightarrow\; x = 48 .
\]
O cone grande tem altura 72. Assim,
\[
V = \tfrac13\cdot 3\,(15^2\cdot 72 - 10^2\cdot 48) = 16\,200 - 4\,800 = 11\,400\ \text{cm}^3,
\]
ou seja, \textbf{11,4~L}.
\end{exemplo}
\end{minipage}

Se preferir, existe uma fórmula pronta para o tronco de cone de raios $R$ e $r$ e
altura $h$, que dá o mesmo resultado:
\[
V_{\text{tronco}} = \frac{\pi h}{3}\left(R^2 + Rr + r^2\right)
\quad\Rightarrow\quad \frac{3\cdot 24}{3}\,(225 + 150 + 100) = 24\cdot 475 = 11\,400\ \text{cm}^3 .
\]

\subsection{Sólidos compostos e deslocamento de líquido}

Silos (cilindro + cone), pílulas (cilindro + duas semiesferas), casquinhas (cone +
hemisfério), peças vazadas (prisma $-$ cilindro)\ldots{} A estratégia é sempre a
mesma: \textbf{decomponha} o objeto em sólidos conhecidos, \textbf{some} o que foi
acrescentado e \textbf{subtraia} o que foi retirado. Cuidado para usar, em cada
parte, a \emph{sua} altura e o \emph{seu} raio.

Outra situação clássica é o \textbf{nível de água que sobe} quando um objeto é
mergulhado. Em um recipiente de paredes verticais, a água “extra” forma um
prisma (ou cilindro) com a mesma base do recipiente e altura igual à subida
$\Delta h$. Como o objeto submerso ocupa exatamente esse espaço:

\begin{formula}[Deslocamento de líquido]
$\displaystyle V_{\text{objeto submerso}} = A_{\text{base do recipiente}}\cdot \Delta h$
\end{formula}

\noindent
\begin{minipage}[c]{0.44\linewidth}
\centering
\begin{tikzpicture}[x={(.78cm,0cm)},y={(.36cm,.26cm)},z={(0cm,.78cm)},line join=round]
  \def\L{5}\def\P{3}\def\A{3}\def\hi{1.5}\def\hf{2.0}
  \draw[v2c12 fina] (0,\P,0)--(0,\P,\A) (0,\P,0)--(\L,\P,0) (0,0,0)--(0,\P,0);
  % pedra
  \fill[ebookMuted!70] (1.3,.6,0) .. controls (1.2,.6,.6) and (1.9,.6,.95) .. (2.4,.6,.7)
    .. controls (2.9,.6,.5) and (2.8,.6,0) .. (2.4,.6,0) -- cycle;
  % água
  \fill[contextoC!22,opacity=.85] (0,0,0)--(\L,0,0)--(\L,0,\hf)--(0,0,\hf)--cycle;
  \fill[contextoC!32,opacity=.85] (\L,0,0)--(\L,\P,0)--(\L,\P,\hf)--(\L,0,\hf)--cycle;
  \fill[contextoC!40,opacity=.8] (0,0,\hf)--(\L,0,\hf)--(\L,\P,\hf)--(0,\P,\hf)--cycle;
  \draw[contextoC,line width=.7pt] (0,0,\hf)--(\L,0,\hf)--(\L,\P,\hf);
  \draw[ebookAcc,line width=.7pt,dash pattern=on 3pt off 2pt] (0,0,\hi)--(\L,0,\hi)--(\L,\P,\hi);
  \draw[v2c12 aresta] (0,0,0)--(\L,0,0)--(\L,\P,0)--(\L,\P,\A)--(0,\P,\A)--(0,0,\A)--cycle
    (0,0,\A)--(\L,0,\A)--(\L,0,0) (\L,0,\A)--(\L,\P,\A);
  \draw[v2c12 cota] (\L+.25,0,\hi)--(\L+.25,0,\hf);
  \node[font=\small,text=ebookAcc,anchor=west] at (\L+.3,0,1.75) {$\Delta h$};
  \node[font=\footnotesize,text=ebookMuted,anchor=east] at (-.1,0,\hi) {nível inicial};
\end{tikzpicture}
\captionof{figure}{A pedra empurra a água para cima: $V_{\text{pedra}} = A_b\cdot\Delta h$.}
\end{minipage}\hfill
\begin{minipage}[c]{0.52\linewidth}
\begin{exemplo}[Volume de uma pedra]
Um aquário tem base retangular de $50\ \text{cm}\times 30\ \text{cm}$. Ao
mergulhar totalmente uma pedra, o nível da água sobe 0,8~cm, sem transbordar.
Qual é o volume da pedra? Se ela tem massa de 3~kg, qual é a sua densidade?
\solucao
$V = 50\cdot 30\cdot 0{,}8 = 1200\ \text{cm}^3$. A densidade é
$\frac{3000\ \text{g}}{1200\ \text{cm}^3} = 2{,}5\ \text{g/cm}^3$.

A forma irregular da pedra não importa: só a área da base do aquário e a
subida do nível.
\end{exemplo}
\end{minipage}

% =====================================================================
\section{Dicas e macetes}

\begin{dica}[Volume ou área? Leia o verbo]
Antes de qualquer conta, decida \emph{o que} está sendo medido. \textbf{Caber,
encher, comportar, armazenar, derreter, despejar, transportar} indicam
\textbf{volume}. \textbf{Pintar, revestir, forrar, embalar, gastar chapa ou
papelão} indicam \textbf{área}. Metade dos erros em questões de sólidos vem de
calcular a grandeza errada com a fórmula certa.
\end{dica}

\begin{dica}[Ache a “base que se repete”]
$V = A_b\cdot h$ vale para \emph{qualquer} prisma ou cilindro, mesmo deitado. A
base é a face que se repete ao longo do sólido, que nem sempre fica embaixo: em
um silo trapezoidal deitado, a base é o trapézio, e a “altura” do prisma é o
comprimento do silo.
\end{dica}

\begin{dica}[Deixe o $\pi$ para o fim]
Ao comparar cilindros, cones e esferas, o $\pi$ aparece em todos os volumes e se
cancela. E, quando o enunciado manda usar $\pi\approx 3$, use exatamente 3 —
as alternativas foram calculadas assim. Se as alternativas trazem o símbolo
$\pi$, não substitua nada.
\end{dica}

\begin{dica}[Nível que sobe ou desce]
Em recipientes de paredes verticais, a variação de volume é sempre
$\Delta V = A_b\cdot\Delta h$. Isso resolve objetos mergulhados, bolinhas
jogadas na água, consumo de reservatórios (“o nível baixou 30~cm”) e vazão
(“subiu 0,6~m em 2~h”).
\end{dica}

\begin{macete}[A escada do volume]
\[
\text{m}^3 \xrightarrow{\ \times 1000\ } \text{dm}^3 = \text{L}
\xrightarrow{\ \times 1000\ } \text{cm}^3 = \text{mL}
\]
Desceu um degrau, multiplique por 1000; subiu, divida por 1000. Um jeito
rápido: converta todas as medidas lineares para \textbf{decímetros} antes de
multiplicar, e o resultado já sai em litros.
\end{macete}

\begin{macete}[$k$, $k^2$, $k^3$]
Se as medidas lineares são multiplicadas por $k$, a área fica multiplicada por
$k^2$ e o volume (ou a massa) por $k^3$. Tabela relâmpago:
\begin{center}
\begin{tabular}{cccc}
\toprule
Medidas & $k$ & Área ($k^2$) & Volume ($k^3$)\\
\midrule
dobram & 2 & $\times 4$ & $\times 8$\\
triplicam & 3 & $\times 9$ & $\times 27$\\
aumentam 50\% & 1,5 & $\times 2{,}25$ & $\times 3{,}375$\\
diminuem 20\% & 0,8 & $\times 0{,}64$ & $\times 0{,}512$\\
caem à metade & 0,5 & $\times 0{,}25$ & $\times 0{,}125$\\
\bottomrule
\end{tabular}
\end{center}
\end{macete}

\begin{macete}[Cilindro: o raio pesa ao quadrado]
Em $V = \pi r^2 h$, dobrar a altura dobra o volume, mas dobrar o raio
\emph{quadruplica} o volume. Para manter o volume quando o raio é multiplicado
por $k$, a altura deve ser dividida por $k^2$. Ex.: raio 2 vezes maior
$\Rightarrow$ altura 4 vezes menor.
\end{macete}

\begin{atencao}[Raio ou diâmetro?]
O ENEM costuma informar o \textbf{diâmetro} (“cano de 10~cm”, “lata de 8~cm de
diâmetro”). As fórmulas usam o \textbf{raio}. Usar o diâmetro no lugar do raio
multiplica a área da base por 4 — e esse valor errado quase sempre está entre
as alternativas.
\end{atencao}

\begin{atencao}[O $\frac13$ é só dos “pontudos”]
Pirâmide e cone levam $\frac13$; prisma e cilindro, não. O hemisfério é
\emph{metade} da esfera: $\frac23\pi r^3$, e não $\frac43\pi r^3$. E a área
lateral do cone usa a \emph{geratriz} $g$, não a altura $h$ — assim como a área
da face de uma pirâmide usa o apótema $m$, não a altura da pirâmide.
\end{atencao}

\begin{atencao}[Metade da altura não é metade do volume]
Em uma taça cônica (vértice para baixo), encher até a metade da altura coloca
só $\left(\frac12\right)^3 = \frac18$ do volume! O líquido forma um cone
\emph{semelhante} ao da taça. Para ter metade do volume, a altura do líquido
deve ser $\sqrt[3]{0{,}5}\approx 0{,}79$ da altura total.
\end{atencao}

\begin{resumo}
\begin{itemize}[leftmargin=1.2em,itemsep=2pt]
  \item \textbf{Prisma e cilindro:} $V = A_b\cdot h$ \quad ($V_{\text{cilindro}} = \pi r^2 h$). Vale também para os oblíquos (Cavalieri).
  \item \textbf{Pirâmide e cone:} $V = \frac13 A_b\cdot h$ \quad ($V_{\text{cone}} = \frac13\pi r^2 h$); \quad $g^2 = h^2 + r^2$.
  \item \textbf{Esfera:} $V = \frac43\pi r^3$, \ $A = 4\pi r^2$; hemisfério: $\frac23\pi r^3$. Cone : esfera : cilindro $= 1:2:3$.
  \item \textbf{Áreas:} cilindro $2\pi r^2 + 2\pi r h$; cone $\pi r^2 + \pi r g$; paralelepípedo $2(ab+ac+bc)$. Some só as faces pedidas.
  \item \textbf{Unidades:} $1\ \text{m}^3 = 1000\ \text{L}$, \ $1\ \text{dm}^3 = 1\ \text{L}$, \ $1\ \text{cm}^3 = 1\ \text{mL}$.
  \item \textbf{Semelhança:} comprimentos $\times k$, áreas $\times k^2$, volumes e massas $\times k^3$.
  \item \textbf{Tronco} $=$ sólido grande $-$ sólido pequeno (use a semelhança para achar a altura que falta).
  \item \textbf{Compostos:} decomponha, some e subtraia. \textbf{Nível da água:} $\Delta V = A_b\cdot\Delta h$.
  \item \textbf{Arredonde com bom senso:} viagens e latas para cima; unidades completas produzidas para baixo.
\end{itemize}
\end{resumo}

% =====================================================================
\secaoenem

% ---------------------------------------------------------------------
\questaoenem{2015-148}
\begin{resolucao}{A}
\passo{1. Entendendo o problema.} Cada contêiner é um paralelepípedo de
$6{,}4\ \text{m}\times 2{,}5\ \text{m}\times 2{,}5\ \text{m}$. Eles devem ocupar a
área de $10\ \text{m}\times 32\ \text{m}$ \emph{sem deixar espaços} e sem
ultrapassá-la. Queremos a altura da pilha.

\passo{2. Quantos contêineres cabem em cada camada.} Para não sobrar espaço, cada
lado da área precisa ser um múltiplo exato da medida do contêiner naquela
direção. Colocando o comprimento de 6,4~m ao longo dos 32~m:
$32 \div 6{,}4 = 5$ (exato) e $10 \div 2{,}5 = 4$ (exato). Logo, cabem
$5\cdot 4 = 20$ contêineres por camada, cobrindo a área inteira. (Na outra
posição, $10\div 6{,}4$ não é inteiro e sobraria espaço.)

\passo{3. Número de camadas e altura.} São $100 \div 20 = 5$ camadas, cada uma
com 2,5~m de altura: $5\cdot 2{,}5 = 12{,}5$~m.

\passo{4. Conferência pelo volume.} Volume total:
$100\cdot(6{,}4\cdot 2{,}5\cdot 2{,}5) = 100\cdot 40 = 4000\ \text{m}^3$. Como
$V = A_b\cdot h$, a altura é $h = 4000 \div (10\cdot 32) = 4000\div 320 = 12{,}5$~m.
Item fácil ($b = 0{,}94$): é o “contar cubinhos em camadas” da teoria.

\resposta{alternativa A — 12,5~m.}
\begin{distratores}
\alt{B} 17,5~m corresponde a 7 camadas; com 20 contêineres por camada, isso
exigiria 140 contêineres, e não 100.
\alt{C} 25,0~m é o dobro do correto: erro de quem considera cada camada com 5~m
de altura (somando as duas medidas de 2,5~m) ou acha que só cabem 10 contêineres
por camada.
\alt{D} Erro de quem gira o contêiner, pondo os 6,4~m ao longo dos 10~m: caberia 1
por fileira e 12 fileiras ($12\cdot 2{,}5 = 30$~m), isto é, 12 por camada e
$\lceil 100/12\rceil = 9$ camadas, ou seja, $9\cdot 2{,}5 = 22{,}5$~m. Mas sobrariam
faixas vazias de 3,6~m e de 2~m, o que a norma proíbe.
\alt{E} 32,5~m seriam 13 camadas (260 contêineres); o valor provavelmente atrai
quem se confunde com a medida de 32~m da área.
\end{distratores}
\end{resolucao}

% ---------------------------------------------------------------------
\questaoenem{2022-169}
\begin{resolucao}{E}
\passo{1. A ideia.} Para um cilindro, $V = A_b\cdot P$ (área da base vezes
profundidade). Vamos escrever cada modelo em função de $V_{\text{I}} = A_b\cdot P$.
O ponto-chave é diferenciar “metade/dobro da \emph{área}” de “metade/dobro do
\emph{raio}”: como $A_b = \pi r^2$, multiplicar o raio por $k$ multiplica a área
da base por $k^2$.

\passo{2. Modelo a modelo.}
\begin{itemize}[nosep,leftmargin=1.2em]
  \item II: $2P\cdot\frac{A_b}{2} = A_b P = V_{\text{I}}$;
  \item III: raio pela metade $\Rightarrow$ área $\frac{A_b}{4}$; $\;2P\cdot\frac{A_b}{4} = \frac12 V_{\text{I}}$;
  \item IV: $\frac{P}{2}\cdot 2A_b = A_b P = V_{\text{I}}$;
  \item V: raio dobrado $\Rightarrow$ área $4A_b$; $\;\frac{P}{2}\cdot 4A_b = 2V_{\text{I}}$.
\end{itemize}

\passo{3. Conclusão.} O modelo V tem o dobro da capacidade do modelo I e é o
maior. Item médio ($b = 1{,}89$): exige perceber que o raio “pesa ao quadrado”.

\resposta{alternativa E — modelo V.}
\begin{distratores}
\alt{A} O modelo I empata com II e IV e perde para V ($2V_{\text{I}}$).
\alt{B} O modelo II tem o mesmo volume do I: dobrar a profundidade e reduzir a
área à metade se compensam. Atrai quem só olha o “dobro da profundidade”.
\alt{C} O modelo III é, na verdade, o \emph{menor} ($\frac12 V_{\text{I}}$). Quem
pensa que reduzir o raio à metade reduz a área só à metade acharia
$2P\cdot\frac{A_b}{2} = V_{\text{I}}$, e ainda assim ele não seria o maior.
\alt{D} O modelo IV tem o mesmo volume do I. Atrai quem confunde “dobro da área”
com “dobro do raio”, que é o que acontece no modelo V.
\end{distratores}
\end{resolucao}

% ---------------------------------------------------------------------
\questaoenem{2024-175}
\begin{resolucao}{D}
\passo{1. O que vira o quê.} Ao enrolar a folha, um dos lados vira a
\textbf{altura} do cilindro e o outro vira o \textbf{comprimento da circunferência
da base}, $2\pi r$. Esse lado \emph{não} é o diâmetro nem o raio.

\passo{2. Embalagem 1} (altura 20~cm, contorno 10~cm):
$2\pi r = 10 \Rightarrow r = \frac{5}{\pi}$. Então
$V_1 = \pi\left(\frac{5}{\pi}\right)^2\cdot 20 = \pi\cdot\frac{25}{\pi^2}\cdot 20 = \frac{500}{\pi}\ \text{cm}^3$.

\passo{3. Embalagem 2} (altura 10~cm, contorno 20~cm):
$2\pi r = 20 \Rightarrow r = \frac{10}{\pi}$. Então
$V_2 = \pi\cdot\frac{100}{\pi^2}\cdot 10 = \frac{1000}{\pi}\ \text{cm}^3$.

\passo{4. Conclusão.} A maior capacidade é a da embalagem 2, com
$\frac{1000}{\pi}\ \text{cm}^3$ (cerca de 318~mL). Note que as duas usam a mesma
quantidade de alumínio na lateral, mas a “mais baixa e larga” guarda o dobro:
o raio pesa ao quadrado no volume. Item médio ($b = 1{,}94$).

\resposta{alternativa D — $\dfrac{1000}{\pi}\ \text{cm}^3$.}
\begin{distratores}
\alt{A} $4000\pi$: usou 20~cm como \emph{raio} e 10~cm como altura
($\pi\cdot 20^2\cdot 10$).
\alt{B} $2000\pi$: usou 10~cm como raio e 20~cm como altura
($\pi\cdot 10^2\cdot 20$).
\alt{C} $\frac{4000}{\pi}$: esqueceu o fator 2 da circunferência, fazendo
$\pi r = 20$, ou seja, $r = \frac{20}{\pi}$ e
$V = \pi\cdot\frac{400}{\pi^2}\cdot 10 = \frac{4000}{\pi}$.
\alt{E} $\frac{500}{\pi}$ é o volume da embalagem 1, a de \emph{menor}
capacidade.
\end{distratores}
\end{resolucao}

% ---------------------------------------------------------------------
\questaoenem{2023-136}
\begin{resolucao}{B}
\passo{1. Plano.} Escultura $=$ cone grande $-$ cone pequeno (do topo) $-$
cilindro (do furo). Depois, massa $=$ densidade $\times$ volume. Use $\pi = 3$.

\passo{2. Altura do cone retirado (semelhança).} O cone grande tem raio 9~cm e
altura 36~cm; o pequeno tem raio 3~cm. Como são semelhantes,
$k = \frac39 = \frac13$, e a altura do cone pequeno é $\frac13\cdot 36 = 12$~cm.
Logo, o tronco tem altura $36 - 12 = 24$~cm.

\passo{3. Volumes.}
\begin{itemize}[nosep,leftmargin=1.2em]
  \item cone grande: $\frac13\cdot 3\cdot 9^2\cdot 36 = 2916\ \text{cm}^3$;
  \item cone pequeno: $\frac13\cdot 3\cdot 3^2\cdot 12 = 108\ \text{cm}^3$;
  \item cilindro do furo (raio 3, atravessa o tronco de base a base, altura 24):
        $3\cdot 3^2\cdot 24 = 648\ \text{cm}^3$.
\end{itemize}
Volume da escultura: $2916 - 108 - 648 = 2160\ \text{cm}^3$.

\passo{4. Massa.} $0{,}6\cdot 2160 = 1296$~g. Item difícil ($b = 2{,}45$): combina
três sólidos, semelhança e densidade.

\resposta{alternativa B — 1\,296,0~g.}
\begin{distratores}
\alt{A} 1\,198,8~g $= 0{,}6\cdot 1998$: o volume do tronco está certo (2\,808),
mas o furo foi calculado com altura 30~cm ($36 - 6$, confundindo o diâmetro de
6~cm com a altura do cone retirado): $2808 - 810 = 1998$.
\alt{C} 1\,360,8~g $= 0{,}6\cdot 2268$: esqueceu de retirar o cone pequeno
($2916 - 648$).
\alt{D} 4\,665,6~g $= 0{,}6\cdot 7776$: esqueceu o $\frac13$ dos cones
($8748 - 324 - 648$).
\alt{E} 4\,860,0~g $= 0{,}6\cdot 8100$: esqueceu o $\frac13$ \emph{e} não retirou
o cone pequeno ($8748 - 648$).
\end{distratores}
\end{resolucao}

% ---------------------------------------------------------------------
\questaoenem{2022-178}
\begin{resolucao}{E}
\passo{1. Semelhança.} Todos os docinhos são esferas, e esferas são semelhantes.
Passar de 2~cm para 4~cm de diâmetro é multiplicar as medidas lineares por
$k = 2$; logo, o volume (e a massa) de cada docinho fica multiplicado por
$k^3 = 8$.

\passo{2. Contando em “docinhos pequenos”.} Cada docinho grande equivale a 8
pequenos. A encomenda equivale a $150\cdot 8 = 1200$ docinhos pequenos.

\passo{3. Porções.} Uma porção rende 50 pequenos:
$1200 \div 50 = 24$ porções. Item difícil ($b = 2{,}57$): a intuição “dobrou o
diâmetro, dobra a massa” é forte, e as alternativas exploram exatamente isso.

\resposta{alternativa E — 24 porções.}
\begin{distratores}
\alt{A} 2: usou apenas a razão entre os diâmetros ($4\div 2$) e ignorou a
quantidade de docinhos.
\alt{B} 3: considerou só a quantidade ($150\div 50$), como se os docinhos não
tivessem mudado de tamanho.
\alt{C} 6: multiplicou por $k = 2$ (proporcionalidade linear): $3\cdot 2$.
\alt{D} 12: multiplicou por $k^2 = 4$ (razão de áreas): $3\cdot 4$.
\end{distratores}
\end{resolucao}

% ---------------------------------------------------------------------
\questaoenem{2010-168}
\begin{resolucao}{B}
\passo{1. Entendendo as taças.} A taça da Figura 1 é um \textbf{hemisfério}
(meia esfera) de raio 3~cm, servida cheia. Na taça cônica, a figura indica que o
champanhe forma um cone com raio 3~cm na superfície e altura $h$, que é o que
queremos.

\passo{2. Volume do hemisfério.}
$V = \frac12\cdot\frac43\pi R^3 = \frac23\pi\cdot 3^3 = 18\pi\ \text{cm}^3$.

\passo{3. Igualando ao cone.}
$\frac13\pi\cdot 3^2\cdot h = 18\pi \Rightarrow 3\pi h = 18\pi \Rightarrow h = 6$~cm.
Repare que o $\pi$ se cancela: nem é preciso substituir valor algum.

Item muito difícil ($b = 3{,}29$), embora as fórmulas sejam dadas: é preciso
perceber que a taça é \emph{meia} esfera e montar a igualdade de volumes; as
alternativas exploram quem usa a esfera inteira ou responde com um volume em vez
de uma altura.

\resposta{alternativa B — 6,00~cm.}
\begin{distratores}
\alt{A} 1,33 $\approx \frac43$: resulta de se perder na simplificação e ficar só
com o coeficiente da fórmula da esfera.
\alt{C} 12,00: usou a esfera \emph{inteira}: $\frac43\pi\cdot 27 = 36\pi = 3\pi h
\Rightarrow h = 12$.
\alt{D} 56,52 $= 18\cdot 3{,}14$: é o \emph{volume} do hemisfério em cm$^3$, e não a
altura pedida.
\alt{E} 113,04 $= 36\cdot 3{,}14$: é o volume da esfera inteira.
\end{distratores}
\end{resolucao}

% =====================================================================
\secaoineditas

% ---------------------------------------------------------------------
\begin{questaoinedita}{1}{12}
Aquários domésticos precisam de um espaço livre entre a superfície da água e a
tampa, para facilitar a troca de gases e evitar que os peixes saltem para fora.
Um aquário com a forma de um paralelepípedo reto-retângulo tem dimensões
internas de 80~cm de comprimento, 40~cm de largura e 50~cm de altura. O
fabricante recomenda que a água fique 5~cm abaixo da borda superior.

\comando{Seguindo a recomendação do fabricante, a quantidade de água, em litro,
que deve ser colocada nesse aquário é}
\begin{alternativas*}[5]
  \item 14,4.
  \item 16,0.
  \item 144,0.
  \item 160,0.
  \item 1\,440,0.
\end{alternativas*}
\end{questaoinedita}
\begin{resolucao}{C}
\passo{1. Altura da água.} A água fica 5~cm abaixo da borda: $50 - 5 = 45$~cm.

\passo{2. Volume.} $V = 80\cdot 40\cdot 45 = 144\,000\ \text{cm}^3$.

\passo{3. Conversão.} Como $1\ \text{L} = 1000\ \text{cm}^3$,
$144\,000 \div 1000 = 144$~L. (Atalho da escada: em decímetros,
$8\cdot 4\cdot 4{,}5 = 144\ \text{dm}^3 = 144$~L.)

\resposta{alternativa C — 144,0~L.}
\begin{distratores}
\alt{A} 14,4: dividiu por \num{10000} na conversão de cm$^3$ para litro.
\alt{B} 16,0: calculou o volume da faixa \emph{vazia} de 5~cm
($80\cdot 40\cdot 5 = 16\,000\ \text{cm}^3$).
\alt{D} 160,0: calculou a capacidade total, com o aquário cheio até a borda.
\alt{E} 1\,440,0: dividiu por 100, e não por 1000, na conversão.
\end{distratores}
\end{resolucao}

% ---------------------------------------------------------------------
\begin{questaoinedita}{2}{8}
Na fábula \emph{O corvo e o jarro}, atribuída a Esopo, um corvo com sede encontra
um jarro com um pouco de água no fundo. Como não alcança a água com o bico, ele
joga pedrinhas dentro do jarro, uma a uma, até o nível subir o suficiente.
Experimentos com corvos reais mostraram que essas aves são, de fato, capazes de
usar essa estratégia.

Em uma montagem que reproduz a fábula, um jarro com a forma de um cilindro
circular reto, de diâmetro interno 8~cm, contém água até a altura de 10~cm. A
ave consegue beber quando o nível da água atinge 16~cm de altura. Cada pedrinha
tem, em média, volume de $6\ \text{cm}^3$ e fica totalmente submersa.

Use 3 como aproximação para $\pi$.

\comando{O número mínimo de pedrinhas que devem ser colocadas no jarro para que
a ave consiga beber a água é}
\begin{alternativas*}[5]
  \item 24.
  \item 48.
  \item 80.
  \item 128.
  \item 192.
\end{alternativas*}
\end{questaoinedita}
\begin{resolucao}{B}
\passo{1. O que precisa ser ocupado.} O nível deve subir de 10~cm para 16~cm:
$\Delta h = 6$~cm. As pedrinhas precisam ocupar o volume de um cilindro com a
base do jarro e altura 6~cm (deslocamento de líquido).

\passo{2. Volume.} Raio $= 8\div 2 = 4$~cm.
$\Delta V = \pi r^2\cdot\Delta h = 3\cdot 4^2\cdot 6 = 288\ \text{cm}^3$.

\passo{3. Número de pedrinhas.} $288\div 6 = 48$. Com 48 pedrinhas o nível chega
exatamente a 16~cm.

\resposta{alternativa B — 48 pedrinhas.}
\begin{distratores}
\alt{A} 24: usou o comprimento da circunferência ($2\pi r$) no lugar da área da
base: $2\cdot 3\cdot 4\cdot 6 = 144$; $144\div 6 = 24$.
\alt{C} 80: calculou o volume da água que já estava no jarro
($3\cdot 16\cdot 10 = 480$; $480\div 6 = 80$).
\alt{D} 128: usou a altura final de 16~cm em vez da variação de 6~cm
($3\cdot 16\cdot 16 = 768$; $768\div 6 = 128$).
\alt{E} 192: usou o diâmetro como raio ($3\cdot 8^2\cdot 6 = 1152$;
$1152\div 6 = 192$).
\end{distratores}
\end{resolucao}

% ---------------------------------------------------------------------
\begin{questaoinedita}{2}{14}
Uma cooperativa rural armazena água em um tanque metálico com a forma de um
cilindro circular reto, com 2~m de raio da base e 5~m de altura, apoiado sobre
uma laje de concreto. Para protegê-lo da ferrugem, a cooperativa vai pintar toda
a parte externa visível do tanque, isto é, a superfície lateral e a tampa; o
fundo, apoiado na laje, não será pintado. A tinta escolhida rende
$8\ \text{m}^2$ por litro, e uma única demão é suficiente.

Use 3 como aproximação para $\pi$.

\comando{A quantidade mínima de tinta, em litro, necessária para essa pintura é}
\begin{alternativas*}[5]
  \item 5,25.
  \item 7,50.
  \item 9,00.
  \item 10,50.
  \item 21,00.
\end{alternativas*}
\end{questaoinedita}
\begin{resolucao}{C}
\passo{1. Quais superfícies?} Pintar é \emph{área}. Entram a lateral e a tampa
(um círculo); o fundo fica de fora.

\passo{2. Áreas.} Lateral: $2\pi r h = 2\cdot 3\cdot 2\cdot 5 = 60\ \text{m}^2$.
Tampa: $\pi r^2 = 3\cdot 2^2 = 12\ \text{m}^2$. Total: $72\ \text{m}^2$.

\passo{3. Tinta.} $72 \div 8 = 9$~L.

\resposta{alternativa C — 9,00~L.}
\begin{distratores}
\alt{A} 5,25: usou $\pi r h$ para a lateral (esqueceu o 2):
$(30 + 12)\div 8 = 5{,}25$.
\alt{B} 7,50: pintou só a lateral ($60\div 8$), esquecendo a tampa.
\alt{D} 10,50: incluiu também o fundo ($60 + 12 + 12 = 84$; $84\div 8$).
\alt{E} 21,00: usou o diâmetro (4~m) como raio:
$(2\cdot 3\cdot 4\cdot 5 + 3\cdot 4^2)\div 8 = 168\div 8 = 21$.
\end{distratores}
\end{resolucao}

% ---------------------------------------------------------------------
\begin{questaoinedita}{3}{8}
O telhado de um quiosque de praia terá o formato da superfície lateral de uma
pirâmide quadrangular regular. A base dessa pirâmide é um quadrado de 6~m de
lado, e o topo do telhado fica 4~m acima do centro desse quadrado, como mostra a
figura.

\begin{center}
\begin{tikzpicture}[x={(.62cm,0cm)},y={(.3cm,.22cm)},z={(0cm,.62cm)},line join=round]
  \fill[ebookAcc!30] (3,3,4)--(0,0,0)--(6,0,0)--cycle;
  \fill[ebookAcc!50] (3,3,4)--(6,0,0)--(6,6,0)--cycle;
  \draw[v2c12 oculta] (0,0,0)--(0,6,0)--(6,6,0) (3,3,4)--(0,6,0);
  \draw[v2c12 oculta] (0,0,0)--(6,6,0) (6,0,0)--(0,6,0);
  \draw[v2c12 aresta] (0,0,0)--(6,0,0)--(6,6,0)--(3,3,4)--cycle (3,3,4)--(6,0,0);
  \draw[ebookPrim,line width=.7pt,dash pattern=on 2pt off 1.5pt] (3,3,4)--(3,3,0);
  \node[font=\small,text=ebookPrim,anchor=west] at (3.05,3,2.1) {4 m};
  \fill[ebookDark] (3,3,0) circle (1pt);
  \draw[v2c12 cota] (0,-.6,0)--(6,-.6,0) node[midway,below,v2c12 rot]{6 m};
  \draw[v2c12 cota] (6.6,0,0)--(6.6,6,0) node[midway,right=2pt,v2c12 rot]{6 m};
\end{tikzpicture}
\end{center}

Para a cobertura, o fornecedor indica o uso de 16 telhas por metro quadrado de
telhado.

\comando{Desconsiderando perdas, o número mínimo de telhas necessárias para
cobrir esse telhado é}
\begin{alternativas*}[5]
  \item 960.
  \item 1\,120.
  \item 1\,344.
  \item 1\,536.
  \item 1\,920.
\end{alternativas*}
\end{questaoinedita}
\begin{resolucao}{A}
\passo{1. Qual área?} O telhado são as 4 faces triangulares (a base quadrada não
é coberta). Cada face é um triângulo de base 6~m; a altura desse triângulo é o
\textbf{apótema da pirâmide} $m$, e não os 4~m de altura da pirâmide — essa é a
informação “escondida”.

\passo{2. Apótema (Pitágoras).} O apótema liga o topo ao ponto médio de um lado da
base. Ele é a hipotenusa de um triângulo retângulo com catetos 4~m (altura) e
3~m (metade do lado): $m = \sqrt{4^2 + 3^2} = \sqrt{25} = 5$~m.

\passo{3. Área do telhado.} $A = 4\cdot\frac{6\cdot 5}{2} = 60\ \text{m}^2$.

\passo{4. Telhas.} $60\cdot 16 = 960$ telhas.

\resposta{alternativa A — 960 telhas.}
\begin{distratores}
\alt{B} 1\,120: usou a aresta lateral ($\sqrt{3^2+3^2+4^2} = \sqrt{34}\approx 5{,}83$~m)
como altura de cada triângulo: $4\cdot\frac{6\cdot 5{,}83}{2}\approx 70\ \text{m}^2$,
e $70\cdot 16 = 1120$.
\alt{C} 1\,344: somou os catetos em vez de usar Pitágoras ($m = 3 + 4 = 7$):
$4\cdot\frac{6\cdot 7}{2} = 84\ \text{m}^2$; $84\cdot 16 = 1344$.
\alt{D} 1\,536: incluiu a base quadrada ($60 + 36 = 96\ \text{m}^2$), que não faz
parte do telhado.
\alt{E} 1\,920: esqueceu de dividir por 2 a área de cada triângulo
($4\cdot 6\cdot 5 = 120\ \text{m}^2$).
\end{distratores}
\end{resolucao}

% ---------------------------------------------------------------------
\begin{questaoinedita}{3}{12}
Uma confeitaria produz bombons esféricos com 3~cm de diâmetro, formados por um
recheio coberto por uma fina camada de chocolate. Em cada bombom são usados
12~g de recheio e 3~g de chocolate de cobertura. Para uma data comemorativa, a
confeitaria lançará uma versão gigante do bombom, também esférica, com 6~cm de
diâmetro e com a mesma espessura de cobertura.

Considere que a massa de recheio é proporcional ao volume do bombom e que a
massa de cobertura, por ser uma camada fina de espessura constante, é
proporcional à área da superfície do bombom.

\comando{A massa total, em grama, de cada bombom gigante será}
\begin{alternativas*}[5]
  \item 30.
  \item 60.
  \item 99.
  \item 108.
  \item 120.
\end{alternativas*}
\end{questaoinedita}
\begin{resolucao}{D}
\passo{1. Razão de semelhança.} As esferas são semelhantes, com
$k = \frac{6}{3} = 2$.

\passo{2. Recheio (volume, $k^3$).} $12\cdot 2^3 = 12\cdot 8 = 96$~g.

\passo{3. Cobertura (área, $k^2$).} $3\cdot 2^2 = 3\cdot 4 = 12$~g.

\passo{4. Total.} $96 + 12 = 108$~g. A dificuldade está em perceber que as duas
partes do bombom crescem em ritmos diferentes: o que “preenche” cresce com o
volume; o que “reveste”, com a área.

\resposta{alternativa D — 108~g.}
\begin{distratores}
\alt{A} 30: multiplicou tudo por $k = 2$ ($15\cdot 2$), como se a massa fosse
proporcional ao diâmetro.
\alt{B} 60: multiplicou tudo por $k^2 = 4$ ($15\cdot 4$), tratando o recheio como
área.
\alt{C} 99: aplicou $k^3$ ao recheio, mas esqueceu de ajustar a cobertura
($96 + 3$).
\alt{E} 120: multiplicou tudo por $k^3 = 8$ ($15\cdot 8$), tratando a cobertura
como volume.
\end{distratores}
\end{resolucao}

% ---------------------------------------------------------------------
\begin{questaoinedita}{4}{14}
Um bar serve licor em taças cuja parte interna tem a forma de um cone circular
reto com o vértice voltado para baixo. O diâmetro da boca da taça mede 8~cm, e a
altura interna, do vértice até a boca, mede 12~cm. Para padronizar o
atendimento, o gerente decidiu gravar em todas as taças um anel indicando o
nível correspondente a uma dose de 81~mL.

\begin{center}
\begin{tikzpicture}[scale=.3]
  \def\ry{.9}
  \fill[ebookAcc!25] (0,0)--(-3,9) arc[start angle=180,end angle=360,x radius=3,y radius=.68] -- cycle;
  \fill[ebookAcc!45] (0,9) ellipse[x radius=3,y radius=.68];
  \draw[v2c12 aresta] (-4,12)--(0,0)--(4,12);
  \draw[v2c12 aresta] (0,12) ellipse[x radius=4,y radius=\ry];
  \draw[ebookAcc!70!black,line width=.9pt] (-3,9) arc[start angle=180,end angle=360,x radius=3,y radius=.68];
  \draw[ebookAcc!70!black,line width=.6pt,dash pattern=on 2pt off 1.5pt] (3,9) arc[start angle=0,end angle=180,x radius=3,y radius=.68];
  \draw[v2c12 aresta] (0,0)--(0,-3);
  \draw[v2c12 aresta] (0,-3) ellipse[x radius=2.6,y radius=.6];
  \draw[v2c12 cota] (5.2,0)--(5.2,12) node[midway,right,v2c12 rot]{12 cm};
  \draw[v2c12 cota] (-4,13.6)--(4,13.6) node[midway,above,v2c12 rot]{8 cm};
  \draw[v2c12 cota] (-4.4,0)--(-4.4,9) node[midway,left,v2c12 rot]{?};
  \node[font=\footnotesize,text=ebookAcc!70!black,anchor=west] at (3.3,8.2) {anel};
\end{tikzpicture}
\end{center}

Considere $1\ \text{cm}^3 = 1\ \text{mL}$ e use 3 como aproximação para $\pi$.

\comando{O anel deve ser gravado a que altura, em centímetro, medida a partir do
vértice da taça?}
\begin{alternativas*}[5]
  \item 3,0.
  \item 5,1.
  \item 6,2.
  \item 7,8.
  \item 9,0.
\end{alternativas*}
\end{questaoinedita}
\begin{resolucao}{E}
\passo{1. Capacidade da taça.} Raio da boca: 4~cm.
$V = \frac13\cdot 3\cdot 4^2\cdot 12 = 192\ \text{cm}^3 = 192$~mL.

\passo{2. O líquido é um cone semelhante.} Com o vértice para baixo, o licor
forma um cone \emph{semelhante} ao da taça. Se a altura do licor é $x$, a razão
de semelhança é $k = \frac{x}{12}$, e a razão entre os volumes é $k^3$:
\[
k^3 = \frac{81}{192} = \frac{27}{64} = \left(\frac34\right)^3
\;\Rightarrow\; k = \frac34 .
\]

\passo{3. Altura do anel.} $x = \frac34\cdot 12 = 9$~cm, medidos a partir do
vértice. Conferência: o licor forma um cone de raio $\frac34\cdot 4 = 3$~cm e
altura 9~cm, com volume $\frac13\cdot 3\cdot 3^2\cdot 9 = 81$~mL.

Item muito difícil: exige calcular a capacidade, perceber a semelhança, trabalhar
com a raiz cúbica (raciocínio inverso) e não confundir “altura a partir do
vértice” com “distância até a borda”.

\resposta{alternativa E — 9,0~cm.}
\begin{distratores}
\alt{A} 3,0: é a distância do anel até a \emph{borda} ($12 - 9$), e não a altura
pedida, medida a partir do vértice.
\alt{B} 5,1: supôs que altura e volume são proporcionais:
$12\cdot\frac{81}{192}\approx 5{,}06$.
\alt{C} 6,2: esqueceu o $\frac13$ do cone (capacidade de 576~mL):
$12\cdot\sqrt[3]{81/576} = 12\cdot\sqrt[3]{9/64}\approx 6{,}24$.
\alt{D} 7,8: usou a raiz quadrada (razão de áreas) em vez da cúbica:
$12\cdot\sqrt{27/64}\approx 7{,}79$.
\end{distratores}
\end{resolucao}
